\documentclass[12pt,a4paper]{article}

% Cần gói hệ thống texlive-lang-other để có mã font T5 tiếng Việt.
\usepackage[T5]{fontenc}
\usepackage[utf8]{inputenc}
\usepackage{lmodern}
\usepackage[vietnamese]{babel}
\usepackage[a4paper,margin=2.5cm]{geometry}
\usepackage{graphicx}
\graphicspath{{figures/}{images/}}
\usepackage{array}
\usepackage{booktabs}
\usepackage{float}
\usepackage{tabularx}
\usepackage[dvipsnames]{xcolor}
\usepackage{tikz}
\usetikzlibrary{arrows.meta,positioning,shapes.geometric,fit}
% Hyperref được nạp sau các gói tạo float để tránh destination trùng lặp.
\usepackage{hyperref}
\linespread{1.25}

\hypersetup{
  colorlinks=true,
  linkcolor=black,
  urlcolor=blue,
  citecolor=black,
  unicode=true,
  pdftitle={Bai thuc hanh 04 - TurtleBot 4: SLAM, Localization va Navigation},
  pdfauthor={Nguyen Van Tong}
}

\setlength{\parindent}{0pt}
\setlength{\parskip}{6pt}
\setlength{\emergencystretch}{3em}

\definecolor{flowblue}{RGB}{232,241,250}
\definecolor{flowgreen}{RGB}{232,246,236}
\definecolor{floworange}{RGB}{252,241,224}
\definecolor{linegray}{RGB}{75,75,75}

\newcommand{\projectpath}{\texttt{hri-assignment-04}}

\tikzset{
  block/.style={draw=linegray, rounded corners, align=center,
                minimum height=9mm, text width=3.2cm, fill=flowblue},
  sensor/.style={block, fill=floworange},
  process/.style={block, fill=flowgreen},
  smallblock/.style={draw=linegray, rounded corners, align=center,
                     minimum height=8mm, text width=2.5cm, fill=flowblue},
  arrow/.style={-{Latex[length=2.2mm]}, thick, draw=linegray},
  dashedarrow/.style={-{Latex[length=2.2mm]}, thick, dashed, draw=linegray}
}

\begin{document}

\pagenumbering{gobble}
\begin{titlepage}
  \centering
  \vspace*{0.8cm}
  {\Large\bfseries ĐẠI HỌC QUỐC GIA HÀ NỘI\par}
  \vspace{0.25cm}
  {\Large\bfseries TRƯỜNG ĐẠI HỌC CÔNG NGHỆ\par}

  \vspace{1.3cm}
  \IfFileExists{Logo.png}{
    \includegraphics[width=4.5cm]{Logo.png}
  }{
    \IfFileExists{images/Logo.png}{
      \includegraphics[width=4.5cm]{images/Logo.png}
    }{
      \fbox{\parbox[c][3.5cm][c]{4.5cm}{\centering Logo\\nhà trường}}
    }
  }

  \vspace{1.3cm}
  {\Huge\bfseries BÀI THỰC HÀNH 04\par}
  \vspace{0.4cm}
  {\LARGE\bfseries TURTLEBOT 4: SLAM, LOCALIZATION\\VÀ NAVIGATION\par}

  \vspace{1.6cm}
  \begin{flushleft}
    \renewcommand{\arraystretch}{1.5}
    \begin{tabular}{@{}>{\bfseries\raggedright\arraybackslash}p{5cm}p{9cm}@{}}
      Họ và tên: & Nguyễn Văn Tổng \\
      Mã sinh viên: & 23020766 \\
      Lớp: & K68E-RE \\
      Giảng viên hướng dẫn: & PGS. TS. Hoàng Văn Xiêm \\
                            & KS. Nguyễn Quốc Bảo \\
    \end{tabular}
  \end{flushleft}
  \vfill
  {\large Hà Nội, 2026\par}
\end{titlepage}

\pagenumbering{roman}
\tableofcontents
\newpage
\listoffigures
\listoftables
\newpage
\pagenumbering{arabic}

\section{Giới thiệu}

Bài thực hành 04 triển khai pipeline điều hướng cho TurtleBot~4 trong mô phỏng.
Robot được khởi chạy trong Gazebo Harmonic, sau đó dữ liệu LiDAR và odometry
được kiểm tra trên ROS~2 và quan sát bằng RViz2. SLAM Toolbox được dùng để khảo
sát môi trường, tạo bản đồ và lưu bản đồ thành hai file \texttt{map.yaml} và
\texttt{map.pgm}.

Sau giai đoạn mapping, SLAM được dừng. Robot được khởi động lại trong cùng môi
trường, bản đồ đã lưu được nạp và AMCL được dùng để định vị. Cuối cùng, Nav2
nhận ba goal khác nhau, lập đường đi và điều khiển robot bằng lệnh vận tốc mà
không cần teleoperation sau khi gửi goal.

Mục tiêu chính của bài là hiểu pipeline điều hướng cơ bản của robot di động
trong ROS~2, thay vì tự xây dựng lại các thuật toán SLAM hoặc navigation.

\section{Môi trường và kiến trúc hệ thống}

\subsection{Môi trường triển khai}

Project được đặt tại \projectpath{}. Bài gốc dùng Ubuntu 22.04 và ROS~2 Humble,
nhưng toàn bộ phần triển khai trong báo cáo này dùng Ubuntu 24.04 và ROS~2
Jazzy. Gazebo Classic không được sử dụng. Các thành phần chính được tóm tắt
trong Bảng~\ref{tab:environment}. Phiên bản package được kiểm tra từ file
\texttt{package.xml} trong bản cài ROS~2 trên máy.

\begin{table}[htbp]
  \centering
  \caption{Môi trường triển khai Bài thực hành 04.}
  \label{tab:environment}
  \renewcommand{\arraystretch}{1.18}
  \begin{tabularx}{\linewidth}{@{}p{0.30\linewidth}X@{}}
    \toprule
    \textbf{Thành phần} & \textbf{Sử dụng} \\
    \midrule
    Hệ điều hành & Ubuntu 24.04.4 LTS \\
    ROS 2 & Jazzy \\
    Simulator & Gazebo Harmonic / Gazebo Sim 8.11.0 \\
    Robot & TurtleBot 4, model từ \texttt{nav2\_minimal\_tb4\_sim} 1.0.1 \\
    Mapping & SLAM Toolbox 2.8.5 \\
    Localization & AMCL và Nav2 Map Server (Nav2 1.3.12) \\
    Navigation & Navigation2 1.3.12 (\texttt{nav2\_bringup}) \\
    Visualization & RViz2 14.1.22 \\
    \bottomrule
  \end{tabularx}
\end{table}

Package tự viết \texttt{hri\_turtlebot4\_navigation} chứa bốn launch file cho
mô phỏng, mapping, localization và navigation. Package cũng chứa cấu hình SLAM,
cấu hình Nav2, world cục bộ, bản đồ đã lưu, cấu hình RViz, node điều khiển WASD và script kiểm tra hệ
thống. Cấu trúc thư mục chính gồm:

\begin{verbatim}
hri_turtlebot4_navigation/
  launch/   simulation, mapping, localization, navigation
  config/   slam_toolbox.yaml, nav2_params.yaml
  worlds/   hri_lab.sdf
  maps/     map.yaml, map.pgm
  rviz/     turtlebot4_navigation.rviz
  scripts/  wasd_teleop.py, verify_system.py
\end{verbatim}

Các package ROS~2 chính và nhiệm vụ của chúng được liệt kê trong
Bảng~\ref{tab:packages}.

\begin{table}[htbp]
  \centering
  \caption{Các package và node chính trong hệ thống.}
  \label{tab:packages}
  \small
  \renewcommand{\arraystretch}{1.15}
  \begin{tabularx}{\linewidth}{@{}p{0.33\linewidth}X@{}}
    \toprule
    \textbf{Package/node} & \textbf{Vai trò} \\
    \midrule
    \texttt{nav2\_minimal\_tb4\_sim}, \texttt{ros\_gz} & Mô phỏng TurtleBot 4, cảm biến, vật lý và cầu nối dữ liệu Gazebo--ROS. \\
    \texttt{robot\_state\_publisher} & Công bố các frame cố định và frame của robot. \\
    \texttt{slam\_toolbox} & Tạo bản đồ từ LiDAR, odometry và TF. \\
    \texttt{map\_server}, \texttt{amcl} & Nạp bản đồ và ước lượng pose của robot trên bản đồ đã biết. \\
    \texttt{planner\_server}, \texttt{controller\_server} & Tạo global path và sinh lệnh điều khiển bám đường. \\
    \texttt{bt\_navigator} & Điều phối quá trình navigation và các hành vi phục hồi. \\
    \texttt{rviz2} & Hiển thị robot, LiDAR, map, TF, costmap, path và công cụ gửi goal. \\
    \bottomrule
  \end{tabularx}
\end{table}

\subsection{Kiến trúc và dữ liệu ROS 2}

Hình~\ref{fig:architecture} cho thấy luồng dữ liệu từ mô phỏng đến mapping,
localization và navigation. SLAM Toolbox chỉ chạy trong giai đoạn tạo bản đồ.
Khi đã có map, Map Server và AMCL thay thế phần SLAM để Nav2 có thể điều hướng
trong môi trường đã biết.

Các tài liệu dùng để đối chiếu cấu hình gồm ROS~2 Jazzy~\cite{ros2-jazzy},
TurtleBot~4 User Manual~\cite{turtlebot4-manual}, repository simulator của
TurtleBot~4~\cite{turtlebot4-simulator}, Gazebo Harmonic~\cite{gazebo-harmonic},
SLAM Toolbox~\cite{slam-toolbox} và Nav2~\cite{nav2}.

\begin{figure}[htbp]
  \centering
  \resizebox{0.94\linewidth}{!}{%
  \begin{tikzpicture}[font=\small, node distance=9mm and 13mm]
    \node[block] (gz) {Gazebo Harmonic};
    \node[block, below=of gz] (tb) {TurtleBot 4};
    \node[sensor, below left=of tb] (scan) {\texttt{/scan}};
    \node[sensor, below right=of tb] (odom) {\texttt{/odom} + TF};
    \node[process, below=17mm of tb] (slam) {SLAM Toolbox};
    \node[block, below=of slam] (map) {\texttt{/map}};
    \node[block, below=of map] (files) {\texttt{map.yaml}\\\texttt{map.pgm}};
    \node[process, below=of files] (amcl) {Map Server + AMCL};
    \node[process, below left=of amcl] (planner) {Global Planner};
    \node[process, below right=of amcl] (controller) {Local Controller};
    \node[block, below=17mm of amcl] (cmd) {\texttt{/cmd\_vel}};

    \draw[arrow] (gz) -- (tb);
    \draw[arrow] (tb) -- (scan);
    \draw[arrow] (tb) -- (odom);
    \draw[arrow] (scan) -- (slam);
    \draw[arrow] (odom) -- (slam);
    \draw[arrow] (slam) -- (map);
    \draw[arrow] (map) -- (files);
    \draw[arrow] (files) -- (amcl);
    \draw[arrow] (amcl) -- (planner);
    \draw[arrow] (amcl) -- (controller);
    \draw[arrow] (planner) -- node[above,font=\scriptsize]{path} (controller);
    \draw[arrow] (controller) -- (cmd);
    \draw[arrow] (cmd.east) -- ++(1.0,0) |- (tb.east);
  \end{tikzpicture}%
  }
  \caption{Luồng dữ liệu chính của hệ thống TurtleBot 4.}
  \label{fig:architecture}
\end{figure}

Các topic đã được kiểm tra trên backend mô phỏng hiện tại được trình bày trong
Bảng~\ref{tab:topics}. Nếu chuyển sang một backend TurtleBot 4 khác thì cần chạy
\texttt{ros2 topic type <topic>} trước khi giữ nguyên cấu hình, đặc biệt đối với
topic vận tốc.

\begin{table}[htbp]
  \centering
  \caption{Các topic chính đã kiểm tra trong hệ thống.}
  \label{tab:topics}
  \small
  \renewcommand{\arraystretch}{1.15}
  \begin{tabularx}{\linewidth}{@{}p{0.19\linewidth}p{0.37\linewidth}X@{}}
    \toprule
    \textbf{Topic} & \textbf{Kiểu dữ liệu} & \textbf{Vai trò} \\
    \midrule
    \texttt{/scan} & \texttt{sensor\_msgs/msg/LaserScan} & Khoảng cách LiDAR theo nhiều góc. \\
    \texttt{/odom} & \texttt{nav\_msgs/msg/Odometry} & Pose và vận tốc cục bộ; giá trị có thể tích lũy sai số. \\
    \texttt{/cmd\_vel} & \texttt{geometry\_msgs/msg/Twist} & Lệnh vận tốc tuyến tính và vận tốc góc. \\
    \texttt{/map} & \texttt{nav\_msgs/msg/OccupancyGrid} & Bản đồ ô lưới do SLAM tạo hoặc Map Server nạp. \\
    \texttt{/tf} & \texttt{tf2\_msgs/msg/TFMessage} & Các phép biến đổi tọa độ thay đổi theo thời gian. \\
    \texttt{/tf\_static} & \texttt{tf2\_msgs/msg/TFMessage} & Các phép biến đổi cố định của robot và sensor. \\
    \texttt{/amcl\_pose} & \texttt{geometry\_msgs/msg/PoseWithCovarianceStamped} & Pose toàn cục do AMCL ước lượng. \\
    \bottomrule
  \end{tabularx}
\end{table}

Trong \texttt{/scan}, mỗi phần tử biểu diễn một khoảng cách ở một góc xác định.
\texttt{/odom} cho biết chuyển động tương đối và thường mượt, nhưng có thể bị
drift theo thời gian. \texttt{/cmd\_vel} dùng \texttt{linear.x} để chạy thẳng
và \texttt{angular.z} để quay. \texttt{/map} là occupancy grid, còn TF liên hệ
dữ liệu được đo trong các hệ tọa độ khác nhau.

Chuỗi TF chính đã kiểm tra là
\texttt{map} $\rightarrow$ \texttt{odom} $\rightarrow$
\texttt{base\_footprint} $\rightarrow$ \texttt{base\_link} $\rightarrow$
\texttt{rplidar\_link}. Hình~\ref{fig:tf-tree} rút gọn chuỗi này để dễ theo
dõi.

\begin{figure}[htbp]
  \centering
  \begin{tikzpicture}[font=\small, node distance=8mm]
    \node[smallblock] (map) {\texttt{map}};
    \node[smallblock, below=of map] (odom) {\texttt{odom}};
    \node[smallblock, below=of odom] (foot) {\texttt{base\_footprint}};
    \node[smallblock, below=of foot] (base) {\texttt{base\_link}};
    \node[smallblock, below=of base] (lidar) {\texttt{rplidar\_link}};
    \draw[arrow] (map) -- node[right,font=\scriptsize]{SLAM hoặc AMCL} (odom);
    \draw[arrow] (odom) -- node[right,font=\scriptsize]{odometry} (foot);
    \draw[arrow] (foot) -- (base);
    \draw[arrow] (base) -- node[right,font=\scriptsize]{sensor TF} (lidar);
  \end{tikzpicture}
  \caption{Chuỗi TF rút gọn của TurtleBot 4.}
  \label{fig:tf-tree}
\end{figure}

Odometry cung cấp chuyển động cục bộ từ \texttt{odom} tới robot. Trong khi
mapping, SLAM Toolbox cung cấp \texttt{map} $\rightarrow$ \texttt{odom}; khi
localization, AMCL cung cấp phép biến đổi này. Nhờ TF, một tia LiDAR đo trong
frame \texttt{rplidar\_link} được đặt đúng vị trí trên bản đồ. Không nên tạo
\texttt{map} $\rightarrow$ \texttt{odom} tĩnh vì cách đó không sửa được sai số
định vị.

\section{Triển khai SLAM, Localization và Nav2}

\subsection{Xây dựng và lưu bản đồ}

Luồng mapping được tóm tắt trong Hình~\ref{fig:slam-flow}. SLAM Toolbox nhận
dữ liệu LiDAR từ \texttt{/scan}, chuyển động từ \texttt{/odom} và quan hệ frame
từ TF để tạo \texttt{/map}. Bản đồ được theo dõi trực tiếp trên RViz2.

\begin{figure}[htbp]
  \centering
  \begin{tikzpicture}[font=\small, node distance=8mm]
    \node[block] (gz) {Gazebo};
    \node[block, below=of gz] (tb) {TurtleBot 4};
    \node[sensor, below=of tb] (data) {\texttt{/scan} + \texttt{/odom} + TF};
    \node[process, below=of data] (slam) {SLAM Toolbox};
    \node[block, below=of slam] (map) {\texttt{/map}};
    \node[block, below=of map] (rviz) {RViz2};
    \draw[arrow] (gz) -- (tb);
    \draw[arrow] (tb) -- (data);
    \draw[arrow] (data) -- (slam);
    \draw[arrow] (slam) -- (map);
    \draw[arrow] (map) -- (rviz);
  \end{tikzpicture}
  \caption{Pipeline xây dựng bản đồ bằng SLAM Toolbox.}
  \label{fig:slam-flow}
\end{figure}

Sau khi mở simulation và launch mapping, robot được điều khiển bằng bàn phím
để đi qua phần lớn khu vực có thể tiếp cận. Robot nên chạy với tốc độ vừa phải,
quay chậm và quay lại một số vị trí cũ để dữ liệu scan có vùng chồng lặp. Quay
quá nhanh hoặc va chạm có thể làm tường trên bản đồ bị lệch hoặc xuất hiện hai
lần. Trong giai đoạn này teleoperation chỉ phục vụ khảo sát môi trường.

Workspace được build trước khi chạy. Mỗi terminal mới cần nạp cả môi trường
ROS~2 lẫn workspace và dùng cùng \texttt{ROS\_DOMAIN\_ID}:

\begin{verbatim}
cd hri-assignment-04
source /opt/ros/jazzy/setup.bash
colcon build --symlink-install
source install/setup.bash
export ROS_DOMAIN_ID=19
\end{verbatim}

Simulation, mapping và teleop được chạy ở ba terminal riêng:

\begin{verbatim}
ros2 launch hri_turtlebot4_navigation simulation.launch.py
ros2 launch hri_turtlebot4_navigation mapping.launch.py
ros2 run hri_turtlebot4_navigation wasd_teleop.py
\end{verbatim}

Các phím \texttt{W}/\texttt{S} điều khiển tiến/lùi;
\texttt{A}/\texttt{D} điều khiển quay trái/phải; Space hoặc \texttt{X} để dừng
và \texttt{Q} để thoát. Robot tự dừng khi không còn nhận phím trong khoảng
0,5~s. Đây là thiết lập của node teleop, không phải thời gian đo từ thực nghiệm.

Khi bản đồ đã đủ rõ, lệnh dưới đây lưu occupancy grid thực tế. Lệnh được chạy
trong thư mục \texttt{src/hri\_turtlebot4\_navigation}:

\begin{verbatim}
ros2 run nav2_map_server map_saver_cli -f maps/map \
  --ros-args -p use_sim_time:=true \
  -p save_map_timeout:=10.0 \
  -p map_subscribe_transient_local:=true
\end{verbatim}

File \texttt{map.pgm} lưu ảnh occupancy grid. File \texttt{map.yaml} lưu đường
dẫn tới ảnh, độ phân giải, gốc bản đồ và ngưỡng phân loại ô trống/ô bị chiếm.
Hai file phải được giữ cùng nhau và được tạo từ chính lần chạy SLAM dùng trong
bài nộp. Sau khi lưu map, workspace được build lại để bản đồ được cập nhật
trong thư mục cài đặt.

\subsection{Localization trên bản đồ đã lưu}

Sau khi lưu map, SLAM Toolbox được dừng và simulation được khởi động lại trong
cùng world. Map Server nạp \texttt{map.yaml} và \texttt{map.pgm}; AMCL kết hợp
bản đồ này với \texttt{/scan} và \texttt{/odom} để ước lượng vị trí robot như
Hình~\ref{fig:localization-flow}.

\begin{figure}[htbp]
  \centering
  \resizebox{0.78\linewidth}{!}{%
  \begin{tikzpicture}[font=\small, node distance=9mm and 14mm]
    \node[block] (files) {\texttt{map.yaml} + \texttt{map.pgm}};
    \node[process, below=of files] (server) {Map Server};
    \node[process, below=of server] (amcl) {AMCL};
    \node[sensor, below left=of amcl] (scan) {\texttt{/scan}};
    \node[sensor, below right=of amcl] (odom) {\texttt{/odom}};
    \node[block, below=17mm of amcl] (pose) {Robot pose trên map};
    \node[block, below=of pose] (rviz) {RViz2};
    \draw[arrow] (files) -- (server);
    \draw[arrow] (server) -- (amcl);
    \draw[arrow] (scan) -- (amcl);
    \draw[arrow] (odom) -- (amcl);
    \draw[arrow] (amcl) -- (pose);
    \draw[arrow] (pose) -- (rviz);
  \end{tikzpicture}%
  }
  \caption{Pipeline localization bằng Map Server và AMCL.}
  \label{fig:localization-flow}
\end{figure}

Launch localization được chạy bằng lệnh:

\begin{verbatim}
ros2 launch hri_turtlebot4_navigation localization.launch.py
\end{verbatim}

Nếu AMCL chưa biết robot đang ở đâu, công cụ \textbf{2D Pose Estimate} trên
RViz2 được dùng để đặt vị trí và hướng gần đúng. Sau đó AMCL so sánh scan hiện
tại với bản đồ và cập nhật pose. SLAM và localization không phải cùng một việc:
SLAM vừa ước lượng chuyển động vừa tạo bản đồ khi chưa biết môi trường, còn AMCL
dùng bản đồ có sẵn để xác định robot đang ở đâu và không tạo bản đồ mới.

\subsection{Điều hướng tự động với Nav2}

Khi Map Server và AMCL đã hoạt động, Nav2 được khởi chạy bằng lệnh:

\begin{verbatim}
ros2 launch hri_turtlebot4_navigation navigation.launch.py
ros2 run hri_turtlebot4_navigation verify_system.py
\end{verbatim}

Người dùng chọn \textbf{Nav2 Goal} trên RViz2, nhấn vào một ô trống trên map và
kéo mũi tên để đặt hướng cuối. Quy trình xử lý goal được minh họa trong
Hình~\ref{fig:nav2-flow}.

\begin{figure}[htbp]
  \centering
  \resizebox{0.94\linewidth}{!}{%
  \begin{tikzpicture}[font=\small, node distance=8mm and 18mm]
    \node[block] (goal) {Nav2 Goal trên RViz2};
    \node[process, below=of goal] (bt) {BT Navigator};
    \node[process, below=of bt] (planner) {Planner Server};
    \node[block, below=of planner] (path) {Global Path};
    \node[process, below=of path] (controller) {Controller Server};
    \node[block, below=of controller] (cmd) {\texttt{/cmd\_vel}};
    \node[block, below=of cmd] (tb) {TurtleBot 4};
    \node[sensor, right=24mm of path] (lidar) {LiDAR};
    \node[process, below=of lidar] (costmap) {Global/Local\\Costmap};
    \draw[arrow] (goal) -- (bt);
    \draw[arrow] (bt) -- (planner);
    \draw[arrow] (planner) -- (path);
    \draw[arrow] (path) -- (controller);
    \draw[arrow] (controller) -- (cmd);
    \draw[arrow] (cmd) -- (tb);
    \draw[arrow] (lidar) -- (costmap);
    \draw[arrow] (costmap) -- (planner);
    \draw[arrow] (costmap) -- (controller);
  \end{tikzpicture}%
  }
  \caption{Pipeline xử lý một navigation goal trong Nav2.}
  \label{fig:nav2-flow}
\end{figure}

BT Navigator điều phối tiến trình navigation. Planner Server tìm đường tổng thể
từ pose hiện tại tới goal trên global costmap. Controller Server bám theo đường
đã tạo, dùng local costmap để phản ứng với vật cản gần robot và sinh lệnh vận
tốc. Costmap biểu diễn vật cản cùng vùng inflation mà robot cần giữ khoảng cách.
Sau khi goal được gửi, robot phải tự chạy bằng Nav2; không dùng bàn phím để hỗ
trợ robot trong lúc này.

\section{Thực nghiệm và kết quả}

Thực nghiệm được chia thành hai giai đoạn. Giai đoạn đầu kiểm tra mô phỏng,
topic và tạo bản đồ bằng teleoperation. Giai đoạn sau khởi động lại cùng world,
nạp map, đặt initial pose bằng \textbf{2D Pose Estimate} rồi chạy Nav2. Các ảnh
dưới đây ghi lại từng bước; kết quả hoàn thành goal cần được đối chiếu với
action status hoặc log Nav2.

Hình~\ref{fig:gazebo} thể hiện robot trong Gazebo.
Hình~\ref{fig:rviz-scan} và Hình~\ref{fig:topic-check} được dùng để đối chiếu
model robot, dữ liệu LiDAR và các topic cơ bản.

\begin{figure}[htbp]
  \centering
  \includegraphics[width=0.85\linewidth,height=7cm,keepaspectratio]{gazebo_turtlebot4.png}
  \caption{TurtleBot 4 trong môi trường mô phỏng Gazebo.}
  \label{fig:gazebo}
\end{figure}

\begin{figure}[htbp]
  \centering
  \includegraphics[width=0.85\linewidth,height=7cm,keepaspectratio]{rviz_robot_scan.png}
  \caption{Robot và dữ liệu LiDAR \texttt{/scan} được quan sát trên RViz2.}
  \label{fig:rviz-scan}
\end{figure}

\begin{figure}[htbp]
  \centering
  \includegraphics[width=0.85\linewidth,height=7cm,keepaspectratio]{topic_verification.png}
  \caption{Kiểm tra các topic cơ bản của TurtleBot 4.}
  \label{fig:topic-check}
\end{figure}


Trong giai đoạn SLAM, bản đồ được theo dõi khi đang xây dựng
(Hình~\ref{fig:slam-progress}) và sau khi khảo sát
(Hình~\ref{fig:completed-map}). Kết quả lưu bản đồ được ghi lại ở
Hình~\ref{fig:map-saved}.

\begin{figure}[htbp]
  \centering
  \includegraphics[width=0.85\linewidth,height=7cm,keepaspectratio]{slam_in_progress.png}
  \caption{Quá trình xây dựng bản đồ bằng SLAM Toolbox.}
  \label{fig:slam-progress}
\end{figure}

\begin{figure}[htbp]
  \centering
  \includegraphics[width=0.85\linewidth,height=7cm,keepaspectratio]{completed_map.png}
  \caption{Bản đồ sau khi robot khảo sát các khu vực chính của môi trường.}
  \label{fig:completed-map}
\end{figure}

\begin{figure}[htbp]
  \centering
  \includegraphics[width=0.85\linewidth,height=7cm,keepaspectratio]{map_saved.png}
  \caption{Bản đồ được lưu sau quá trình SLAM.}
  \label{fig:map-saved}
\end{figure}


Khi localization, scan cần khớp tương đối với tường trên bản đồ và cụm hạt AMCL
cần tập trung quanh robot (Hình~\ref{fig:amcl}). Sau đó Nav2 dùng costmap để
lập đường đi như Hình~\ref{fig:nav2-path}.

\begin{figure}[htbp]
  \centering
  \includegraphics[width=0.85\linewidth,height=7cm,keepaspectratio]{amcl_localization.png}
  \caption{Localization của TurtleBot 4 trên bản đồ đã lưu.}
  \label{fig:amcl}
\end{figure}

\begin{figure}[htbp]
  \centering
  \includegraphics[width=0.85\linewidth,height=7cm,keepaspectratio]{nav2_path_costmap.png}
  \caption{Nav2 lập đường đi từ vị trí hiện tại tới navigation goal.}
  \label{fig:nav2-path}
\end{figure}


Ba goal được chọn lần lượt tại các vùng trống khác nhau trên RViz2
(Hình~\ref{fig:goal-1}--\ref{fig:goal-3}). Teleop cần được thoát trước khi gửi
goal để Nav2 điều khiển robot. Ảnh đường đi hoặc vị trí robot không tự xác nhận goal
đã hoàn thành; trạng thái \texttt{SUCCEEDED} của action hoặc log
\texttt{Goal succeeded} được dùng để xác nhận kết quả.

\begin{figure}[htbp]
  \centering
  \includegraphics[width=0.85\linewidth,height=7cm,keepaspectratio]{goal_1.png}
  \caption{Navigation Goal 1 trên bản đồ.}
  \label{fig:goal-1}
\end{figure}

\begin{figure}[htbp]
  \centering
  \includegraphics[width=0.85\linewidth,height=7cm,keepaspectratio]{goal_2.png}
  \caption{Navigation Goal 2 trên bản đồ.}
  \label{fig:goal-2}
\end{figure}

\begin{figure}[htbp]
  \centering
  \includegraphics[width=0.85\linewidth,height=7cm,keepaspectratio]{goal_3.png}
  \caption{Navigation Goal 3 trên bản đồ.}
  \label{fig:goal-3}
\end{figure}

\begin{figure}[htbp]
  \centering
  \includegraphics[width=0.85\linewidth,height=7cm,keepaspectratio]{navigation_result.png}
  \caption{Trạng thái navigation goal được ghi nhận từ Nav2.}
  \label{fig:navigation-result}
\end{figure}

Bảng~\ref{tab:results} phân biệt các bước đã có bằng chứng kiểm tra với ba
goal trong phiên RViz cần xác nhận riêng. Kết quả kiểm thử goal tự động trước đó
không thay thế cho ba goal gửi từ RViz2 trong bài nộp.

\begin{table}[htbp]
  \centering
  \caption{Tổng hợp kết quả và bằng chứng thực nghiệm.}
  \label{tab:results}
  \small
  \renewcommand{\arraystretch}{1.14}
  \begin{tabularx}{\linewidth}{@{}p{0.35\linewidth}p{0.22\linewidth}X@{}}
    \toprule
    \textbf{Thử nghiệm} & \textbf{Kết quả} & \textbf{Ghi chú} \\
    \midrule
    Gazebo khởi chạy TurtleBot 4 & Đã kiểm tra & Mô phỏng robot và world cục bộ. \\
    \texttt{/scan} hoạt động & Đã kiểm tra & Kiểu \texttt{LaserScan}; được hiển thị trên RViz2. \\
    \texttt{/odom} hoạt động & Đã kiểm tra & Kiểu \texttt{Odometry}; có dữ liệu chuyển động. \\
    Robot nhận lệnh vận tốc & Đã kiểm tra & Backend nhận \texttt{Twist} tại \texttt{/cmd\_vel}. \\
    SLAM tạo được bản đồ & Đã kiểm tra & Theo dõi \texttt{/map} và TF trong giai đoạn mapping. \\
    Lưu \texttt{map.yaml}, \texttt{map.pgm} & Đã kiểm tra & Map Saver tạo hai file trong \texttt{maps/}. \\
    Load lại bản đồ & Đã kiểm tra & Map Server hoạt động ở trạng thái \texttt{active}. \\
    Localization hoạt động & Đã kiểm tra & AMCL \texttt{active}; scan khớp gần với tường trên map. \\
    Navigation Goal 1 trên RViz2 & \textit{[Xác nhận từ log]} & Đối chiếu Hình~\ref{fig:goal-1} với action status của goal này. \\
    Navigation Goal 2 trên RViz2 & \textit{[Xác nhận từ log]} & Đối chiếu Hình~\ref{fig:goal-2} với action status của goal này. \\
    Navigation Goal 3 trên RViz2 & \textit{[Xác nhận từ log]} & Đối chiếu Hình~\ref{fig:goal-3} với action status của goal này. \\
    \bottomrule
  \end{tabularx}
\end{table}

\section{Một số vấn đề gặp phải}

Hai vấn đề quan sát được trong quá trình chạy được ghi lại ở
Bảng~\ref{tab:issues}.

\begin{table}[htbp]
  \centering
  \caption{Các vấn đề đã gặp và hướng xử lý.}
  \label{tab:issues}
  \small
  \renewcommand{\arraystretch}{1.17}
  \begin{tabularx}{\linewidth}{@{}p{0.27\linewidth}p{0.30\linewidth}X@{}}
    \toprule
    \textbf{Vấn đề} & \textbf{Kiểm tra/nguyên nhân} & \textbf{Cách xử lý} \\
    \midrule
    Scan lệch khỏi tường khi localization
      & Ảnh RViz2 ban đầu cho thấy scan đỏ lệch rõ khỏi tường. Pose ban đầu chưa phù hợp là một khả năng cần kiểm tra.
      & Dùng \textbf{2D Pose Estimate} đặt lại vị trí và hướng, di chuyển chậm để AMCL cập nhật rồi đối chiếu scan với map. Ảnh sau cho thấy scan khớp sát hơn. \\
    Gửi goal nhưng robot chưa di chuyển; không thấy costmap
      & Map Server và AMCL đã \texttt{active}, nhưng chưa có các node planner, controller và BT Navigator. Hai topic costmap có publisher count bằng 0.
      & Khởi chạy \texttt{navigation.launch.py}; kiểm tra các node Nav2 ở trạng thái \texttt{active} và costmap có publisher trước khi gửi lại goal. \\
    \bottomrule
  \end{tabularx}
\end{table}

Một topic xuất hiện trong \texttt{ros2 topic list} chưa đủ chứng minh nó có
dữ liệu: RViz có thể đã tạo subscriber cho topic đó. Vì vậy cần kiểm tra
publisher bằng \texttt{ros2 topic info}, trạng thái lifecycle và TF trước khi
điều chỉnh tham số SLAM hoặc Nav2.

\section{Kết luận}

Bài thực hành xây dựng một quy trình gồm mô phỏng TurtleBot 4, kiểm tra LiDAR và
odometry, tạo rồi lưu bản đồ bằng SLAM Toolbox, nạp lại bản đồ để localization
và chạy Nav2 với ba navigation goal. Các bước mô phỏng, mapping, lưu map và localization đã được kiểm tra.
Ba goal gửi từ RViz2 cần được đối chiếu với log hoặc action status của từng
lần chạy để hoàn tất Bảng~\ref{tab:results}.

Qua quá trình triển khai, em hiểu rõ hơn vai trò của ROS topic và TF trong việc nối
dữ liệu cảm biến với robot và bản đồ. Đồng thời, bài phân biệt rõ mapping với
localization, cũng như luồng xử lý của Nav2 từ goal, global path và costmap tới
lệnh \texttt{/cmd\_vel} làm robot tự di chuyển.

\section*{Liên kết sản phẩm}
\addcontentsline{toc}{section}{Liên kết sản phẩm}

\begin{itemize}
  \item GitHub repository Bài 04:
    \href{https://github.com/VTongg159/hri-assignment-04}
    {\texttt{github.com/VTongg159/hri-assignment-04}}
  \item Video demo: \textit{[Điền link Google Drive]}
\end{itemize}

\begin{thebibliography}{9}
  \bibitem{ros2-jazzy}
  Open Robotics, \emph{ROS 2 Jazzy documentation},
  \url{https://docs.ros.org/en/jazzy/}.

  \bibitem{turtlebot4-manual}
  Clearpath Robotics, \emph{TurtleBot 4 User Manual},
  \url{https://turtlebot.github.io/turtlebot4-user-manual/}.

  \bibitem{turtlebot4-simulator}
  TurtleBot, \emph{TurtleBot 4 Simulator repository},
  \url{https://github.com/turtlebot/turtlebot4_simulator}.

  \bibitem{slam-toolbox}
  Steve Macenski et al., \emph{SLAM Toolbox documentation},
  \url{https://docs.ros.org/en/ros2_packages/jazzy/api/slam_toolbox/}.

  \bibitem{nav2}
  Open Navigation LLC, \emph{Navigation2 documentation},
  \url{https://docs.nav2.org/}.

  \bibitem{gazebo-harmonic}
  Open Robotics, \emph{Gazebo Harmonic documentation},
  \url{https://gazebosim.org/docs/harmonic/}.
\end{thebibliography}

\end{document}
